
\documentclass[11pt]{article}
\usepackage[margin=0.5in]{geometry}

% math packages
\usepackage{amssymb}
\usepackage{amsmath}
\usepackage{amsthm}
\usepackage{mathtools}

% quantum packages
\usepackage{braket}
\usepackage{quantikz}

% lists
\usepackage{enumitem}

% color + links
\usepackage{xcolor}
\usepackage{hyperref}
\hypersetup{
  colorlinks=true,
  linkcolor=blue,
  urlcolor=blue,
  filecolor=blue
}

% line spacing & formatting
\usepackage{parskip}
\usepackage{soul}

% custom commands
\newcommand{\invsq}{\frac{1}{\sqrt{2}}}
\newcommand{\sinth}{\sin\theta}
\newcommand{\costh}{\cos\theta}
\newcommand{\ztimesn}{0^{\otimes n}}
\newcommand{\ketphi}{\ket{\phi}}
\newcommand{\ketpsi}{\ket{\psi}}
\DeclarePairedDelimiter\abs{\lvert}{\rvert}

% title stuff
\title{CSCI 5244: Homework Set 4}
\date{}
\author{Rey Stone}


% ----------- DOCUMENT STARTS HERE ------------

\begin{document}
\maketitle

\raggedright
\begin{enumerate}
	\item \textbf{[Phase Flip Codes]} Recall the phase-flip error $\begin{cases}\ket{0} \rightarrow \ket{0} \\ \ket{1} \rightarrow -\ket{1} \end{cases}$, and the states $\ket{+} = \invsq (\ket{0} + \ket{1})$ and $\ket{1} = \invsq (\ket{0} - \ket{1})$. The following $[3,1]$ codes can be used to detect and correct the phase-flip error:

	      \begin{equation*}
		      \begin{cases}
			      \ket{0} \rightarrow \ket{\bar{0}} = \ket{+++} \\
			      \ket{1} \rightarrow \ket{\bar{1}} = \ket{---}
		      \end{cases}
	      \end{equation*}

	      \begin{enumerate}
		      \item Draw the quantum circuit for encoding the above phase-flip codes.

		            \textbf{Begin solution}

		            \begin{quantikz}[wire types={q,q,q,n,n}]
			            \lstick{$\ket{\phi}$} & & \ctrl{1} & \ctrl{2} & \gate{H} & \\
			            \lstick{$\ket{0}$} & & \targ{} & & \gate{H} & \\
			            \lstick{$\ket{0}$} & & & \targ{} & \gate{H} &
		            \end{quantikz}

		      \item Give the stabilizers that can be used to detect the phase-flip error.

		            \textbf{Begin solution}

		            Since this is a phase-flip error, we use the Pauli \textbf{X} matrices to flip $\ket{+}$ to $\ket{-}$.

		            The stabilizers would be

		            \begin{align*}
			            M_1 & = X_1 \otimes X_2 \otimes I \\
			            M_2 & = I \otimes X_2 \otimes X_3
		            \end{align*}


		      \item Continuing on the encoding circuit from (a), draw the circuit for error detection and correction.

		            \textbf{Begin solution}

		            \begin{quantikz}[wire types={q,q,q,n,n}]
			            \lstick{$\ket{\phi}$} & & \ctrl{1} & \ctrl{2} & \gate{H} & \gate[3]{E} & \gate{H} & & \ctrl{3} & & & & & \gate{X} & & & \gate{H} \\
			            \lstick{$\ket{0}$} & & \targ{} & & \gate{H} & & \gate{H} & & & \ctrl{2} & \ctrl{3} & & & & \gate{X} & & \gate{H} \\
			            \lstick{$\ket{0}$} & & & \targ{} & \gate{H} & & \gate{H} & & & & & \ctrl{2} & & & & \gate{X} & \gate{H} \\
			            &&&&&&&\lstick{$\ket{0}$} & \targ{}\setwiretype{q} & \targ{} & & & \meter[2]{} & \ctrl[vertical wire=c]{-3}\setwiretype{c} & \ctrl[vertical wire=c]{-2} & \octrl[vertical wire=c]{-1} \\
			            &&&&&&&\lstick{$\ket{0}$} &\setwiretype{q} & & \targ{} & \targ{} & & \octrl[vertical wire=c]{-1}\setwiretype{c} & \ctrl[vertical wire=c]{-1} & \ctrl[vertical wire=c]{-1} \\
		            \end{quantikz}

	      \end{enumerate}

	\item Recall Pauli X,Y,Z matrices: $\sigma_1 =
		      \begin{pmatrix}
			      0 & 1 \\ 1 & 0
		      \end{pmatrix},
		      \sigma_2 =
		      \begin{pmatrix}
			      0 & -i \\ i & 0
		      \end{pmatrix},
		      \sigma_3 =
		      \begin{pmatrix}
			      1 & 0 \\ 0 & -1
		      \end{pmatrix}
	      $

	      \begin{enumerate}
		      \item Show  the following anti-commutation relationships:

		            \begin{equation*}
			            \{\sigma_i, \sigma_j\} = 2 \delta_{ij}I
		            \end{equation*}

		            where $\delta_{ij}$ is the Kronecker delta.

		            \textbf{Begin solution}

		            An anti-commutator is a bilinear map defined as $\{X,Y\} = XY + YX$ and has the following properties

		            \begin{itemize}
			            \item $\{A,B\}^{\dagger} = \{B^\dagger, A^\dagger\}$
			            \item $i\{A,B\}^\dagger = -i{A,B}$
		            \end{itemize}

		            (Taken from \href{https://publish.obsidian.md/myquantumwell/Quantum+Mechanics/Anti-commutators+in+quantum+mechanics}{The Quantum Well})

		            The Kronecker delta is defined as: $\delta_{ij} = 1$ if $i = j $ else $0$ if $i \ne j$.

		            To check these claims we can first check when $i = j = 1$ for $\{\sigma_1, \sigma_1\}$

		            \begin{align*}
			            \sigma_1 \sigma_1                     & =
			            \begin{bmatrix} 0 & 1 \\ 1 & 0 \end{bmatrix}
			            \begin{bmatrix} 0 & 1 \\ 1 & 0 \end{bmatrix} =
			            \begin{bmatrix}
				            1 & 0 \\ 0 & 1
			            \end{bmatrix}                             \\
			            \sigma_1 \sigma_1 + \sigma_1 \sigma_1 & =
			            \begin{bmatrix}
				            1 & 0 \\ 0 & 1
			            \end{bmatrix} +
			            \begin{bmatrix}
				            1 & 0 \\ 0 & 1
			            \end{bmatrix}                             \\
			                                                  & = 2
			            \begin{bmatrix}
				            1 & 0 \\ 0 & 1
			            \end{bmatrix} = 2I
		            \end{align*}

		            Next we can check for $\sigma_1 \sigma_2$ which is

		            \begin{align*}
			            \sigma_1 \sigma_2                     & =
			            \begin{bmatrix} 0 & 1 \\ 1 & 0 \end{bmatrix}
			            \begin{bmatrix} 0 & -i \\ i & 0 \end{bmatrix} =
			            \begin{bmatrix}
				            i & 0 \\ 0 & -i
			            \end{bmatrix}                           \\
			            \sigma_2 \sigma_1                     & =
			            \begin{bmatrix} 0 & -i \\ i & 0 \end{bmatrix}
			            \begin{bmatrix} 0 & 1 \\ 1 & 0 \end{bmatrix} =
			            \begin{bmatrix}
				            -i & 0 \\ 0 & i
			            \end{bmatrix}                           \\
			            \sigma_1 \sigma_2 + \sigma_1 \sigma_1 & =
			            \begin{bmatrix}
				            i & 0 \\ 0 & -i
			            \end{bmatrix} +
			            \begin{bmatrix}
				            -i & 0 \\ 0 & i
			            \end{bmatrix} = 0
		            \end{align*}

		            For $\sigma_2$ and $\sigma_3$

		            \begin{align*}
			            \sigma_2 \sigma_3                     & =
			            \begin{bmatrix} 0 & -i \\ i & 0 \end{bmatrix}
			            \begin{bmatrix} 1 & 0 \\ 0 & -1 \end{bmatrix} =
			            \begin{bmatrix}
				            0 & i \\ i & 0
			            \end{bmatrix}                           \\
			            \sigma_3 \sigma_2                     & =
			            \begin{bmatrix} 1 & 0 \\ 0 & -1 \end{bmatrix}
			            \begin{bmatrix} 0 & -i \\ i & 0 \end{bmatrix} =
			            \begin{bmatrix}
				            0 & -i \\ -i & 0
			            \end{bmatrix}                          \\
			            \sigma_2 \sigma_3 + \sigma_3 \sigma_2 & =
			            \begin{bmatrix}
				            0 & i \\ i & 0
			            \end{bmatrix} +
			            \begin{bmatrix}
				            0 & -i \\ -i & 0
			            \end{bmatrix} = 0
		            \end{align*}

		            Finally for $\sigma_1$ and $\sigma_3$

		            \begin{align*}
			            \sigma_1 \sigma_3                     & =
			            \begin{bmatrix} 0 & 1 \\ 1 & 0 \end{bmatrix}
			            \begin{bmatrix} 1 & 0 \\ 0 & -1 \end{bmatrix} =
			            \begin{bmatrix}
				            0 & -1 \\ 1 & 0
			            \end{bmatrix}                           \\
			            \sigma_3 \sigma_1                     & =
			            \begin{bmatrix} 1 & 0 \\ 0 & -1 \end{bmatrix} =
			            \begin{bmatrix} 0 & 1 \\ 1 & 0 \end{bmatrix}
			            \begin{bmatrix}
				            0 & 1 \\ -1 & 0
			            \end{bmatrix}                           \\
			            \sigma_1 \sigma_3 + \sigma_3 \sigma_1 & =
			            \begin{bmatrix}
				            0 & -1 \\ 1 & 0
			            \end{bmatrix} +
			            \begin{bmatrix}
				            0 & 1 \\ -1 & 0
			            \end{bmatrix} = 0
		            \end{align*}


		      \item Show the following commutation relationships:

		            \begin{equation*}
			            [\sigma_j, \sigma_k] = 2i\epsilon_{jkl} \sigma_l
		            \end{equation*}

		            where $\epsilon_{jkl}$ is the Levi-Civita symbol.

		            \textbf{Begin solution}

		            First, we can exapand $[\sigma_j, \sigma_k] = \sigma_j \sigma_k - \sigma_k \sigma_j$.

		            To check these claims we can first check when $j = 1, k =2$. This gives us

		            \begin{equation*}
			            2i(\epsilon_{121} \sigma_1 + \epsilon_{122} \sigma_2 + \epsilon_{123} \sigma_3) = 2i\sigma_3
		            \end{equation*}

		            \begin{align*}
			            \sigma_1 \sigma_2                     & =
			            \begin{bmatrix}
				            0 & 1 \\ 1 & 0
			            \end{bmatrix}
			            \begin{bmatrix}
				            0 & -i \\ i & 0
			            \end{bmatrix} =
			            \begin{bmatrix}
				            i & 0 \\ 0 & -i
			            \end{bmatrix}                           \\
			            \sigma_2 \sigma_1                     & =
			            \begin{bmatrix}
				            0 & -i \\ i & 0
			            \end{bmatrix}
			            \begin{bmatrix}
				            0 & 1 \\ 1 & 0
			            \end{bmatrix} =
			            \begin{bmatrix}
				            -i & 0 \\ 0 & i
			            \end{bmatrix}                           \\
			            \sigma_1 \sigma_2 - \sigma_2 \sigma_1 & =
			            \begin{bmatrix}
				            i & 0 \\ 0 & -i
			            \end{bmatrix} -
			            \begin{bmatrix}
				            -i & 0 \\ 0 & i
			            \end{bmatrix} =
			            2i \begin{bmatrix}
				               1 & 0 \\ 0 & -1
			               \end{bmatrix} = 2i\sigma_3
		            \end{align*}

		            Check $j = 2, k = 3$

		            \begin{equation*}
			            2i(\epsilon_{231} \sigma_1 + \epsilon_{232} \sigma_2 + \epsilon_{233} \sigma_3) = 2i\sigma_1
		            \end{equation*}

		            \begin{align*}
			            \sigma_2 \sigma_3                     & =
			            \begin{bmatrix}
				            0 & -i \\ i & 0
			            \end{bmatrix}
			            \begin{bmatrix}
				            1 & 0 \\ 0 & -1
			            \end{bmatrix} =
			            \begin{bmatrix}
				            0 & i \\ i & 0
			            \end{bmatrix}                           \\
			            \sigma_3 \sigma_2                     & =
			            \begin{bmatrix}
				            1 & 0 \\ 0 & -1
			            \end{bmatrix}
			            \begin{bmatrix}
				            0 & 1 \\ 1 & 0
			            \end{bmatrix} =
			            \begin{bmatrix}
				            0 & -i \\ -i & 0
			            \end{bmatrix}                          \\
			            \sigma_2 \sigma_3 - \sigma_3 \sigma_2 & =
			            \begin{bmatrix}
				            0 & i \\ i & 0
			            \end{bmatrix} -
			            \begin{bmatrix}
				            0 & -i \\ -i & 0
			            \end{bmatrix} =
			            2i \begin{bmatrix}
				               0 & 1 \\ 1 & 0
			               \end{bmatrix} = 2i\sigma_1
		            \end{align*}

		            Check $j = 3, k = 1$

		            \begin{equation*}
			            2i(\epsilon_{311} \sigma_1 + \epsilon_{312} \sigma_2 + \epsilon_{313} \sigma_3) = 2i\sigma_2
		            \end{equation*}

		            \begin{align*}
			            \sigma_3 \sigma_1 & =
			            \begin{bmatrix}
				            1 & 0 \\ 0 & -1
			            \end{bmatrix}
			            \begin{bmatrix}
				            0 & 1 \\ 1 & 0
			            \end{bmatrix} =
			            \begin{bmatrix}
				            0 & 1 \\ -1 & 0
			            \end{bmatrix}       \\
			            \sigma_1 \sigma_3 & =
			            \begin{bmatrix}
				            0 & 1 \\ 1 & 0
			            \end{bmatrix}
			            \begin{bmatrix}
				            1 & 0 \\ 0 & -1
			            \end{bmatrix} =
			            \begin{bmatrix}
				            0 & -1 \\ 1 & 0
			            \end{bmatrix}       \\
			            \sigma_3 \sigma_1 - \sigma_1 \sigma_3 & =
			            \begin{bmatrix}
				            0 & 1 \\ -1 & 0
			            \end{bmatrix} -
			            \begin{bmatrix}
				            0 & -1 \\ 1 & 0
			            \end{bmatrix} =
			            2 \begin{bmatrix}
                      0 & 1 \\ -1 & 0
			               \end{bmatrix} = 2i\sigma_2
		            \end{align*}

                For $j = k = 1$ we get

                \begin{equation*}
			            2i(\epsilon_{111} \sigma_1 + \epsilon_{112} \sigma_2 + \epsilon_{113} \sigma_3) = 0
                \end{equation*}

                \begin{align*}
                  \sigma_1 \sigma_1 - \sigma_1 \sigma_1 =
                  \begin{bmatrix}
                    0 & 1 \\ 1 & 0
                  \end{bmatrix}
                  \begin{bmatrix}
                    0 & 1 \\ 1 & 0
                  \end{bmatrix} -
                  \begin{bmatrix}
                    0 & 1 \\ 1 & 0
                  \end{bmatrix}
                  \begin{bmatrix}
                    0 & 1 \\ 1 & 0
                  \end{bmatrix} = 0
                \end{align*}

	      \end{enumerate}

\end{enumerate}

% ----------- DOCUMENT ENDS HERE ------------
\end{document}
